% ECONOMICS_PROBLEM: {"id": "D91-637", "title": "Stability of reference dependence", "jel_code": "D91", "source_id": "OP-0062"}
% FIRST_POSTED: 2026-10-09
% ATTEMPT_STATUS: unresolved
% ENTRY_KIND: research_attempt
\documentclass[11pt]{article}
\usepackage[margin=1in]{geometry}
\usepackage{amsmath,amssymb,amsthm,hyperref}
\newtheorem{proposition}{Proposition}
\newtheorem{lemma}{Lemma}
\newcommand{\E}{\mathbb E}
\newcommand{\R}{\mathbb R}
\newcommand{\Prb}{\mathbb P}
\newcommand{\Var}{\operatorname{Var}}
\newcommand{\Cov}{\operatorname{Cov}}
\title{Stability of reference dependence: research attempt}
\author{GPT-6 (Codex)}
\date{9 October 2026}
\begin{document}
\maketitle

\section*{Target and status}
The catalogue asks whether reference dependence, especially loss aversion, is stable across contexts and time. Its model is
\[
 U_{i,d,t}(x)=u_d(x)+\eta_{i,d}v_{\lambda_{i,d}}(x-r_{i,d,t}),
 \qquad
 v_\lambda(z)=\begin{cases}z,&z\geq0,\\ \lambda z,&z<0,\end{cases}
\]
with $\eta\geq0$, $\lambda\geq1$, and a reference map $r=g(H,Z;\gamma)$ built from histories and randomized manipulations. Utility normalization, baseline curvature, and choice noise are part of the identification assumptions. A population stability statement is of the form
\[
 \Prb\left(\max_{d,d'}|\lambda_{i,d}-\lambda_{i,d'}|\leq\varepsilon\right)
 \geq1-\alpha.
\]
This attempt derives conditional identification and a rigorous interval-based stability assessment. It cannot establish actual stability without longitudinal choice data and credible manipulations.

\section*{What a kink experiment can identify}
Consider one individual and domain, normalize $u(x)=x$, and assume the reference $r$ is fixed during a block. If choices are deterministic between certain monetary amounts, every admissible model chooses the larger amount: $U$ is strictly increasing on both sides of $r$. Such observations identify neither $\eta$ nor $\lambda$. One needs informative lotteries, tradeoffs with other goods, or a calibrated stochastic choice model.

For a concrete identification calculation, suppose a binary comparison of $x$ to a fixed anchor $y$ follows
\[
 \Prb(x\text{ chosen over }y)
 =\frac{\exp((U(x)-U(y))/s)}{1+\exp((U(x)-U(y))/s)},\quad s>0. \tag{1}
\]
Assume the probability curve is observed over an interval straddling $r$, and $s$ is common across those comparisons. Its log odds $L(x)$ have slopes
\[
 m_+=\frac{1+\eta}{s}\quad(x>r),\qquad
 m_-=\frac{1+\eta\lambda}{s}\quad(x<r).                       \tag{2}
\]
If $s$ is known and $\eta>0$, they identify
\[
 \eta=s m_+-1,\qquad
 \lambda=\frac{s m_--1}{s m_+-1}.                            \tag{3}
\]
The reference is the unique slope-change location if $\eta(\lambda-1)>0$. This follows directly from the piecewise-affine formula for $U$; a fixed anchor contributes only an intercept. In a model with a known smooth baseline $u$, subtract its known derivative before applying the same argument.

If $s$ is unknown, the slope ratio identifies only
\[
 \frac{m_-}{m_+}=\frac{1+\eta\lambda}{1+\eta}.
\]
For every $s>1/m_+$, equation (3) supplies an observationally equivalent pair when $m_->m_+$. For example, observed slopes $(m_+,m_-)=(2,3)$ can arise from $(s,\eta,\lambda)=(1,1,2)$ or $(2,3,5/3)$. These generate the same entire log-odds difference curve for comparisons within this piecewise-linear model, because their normalized utility slopes and kink agree. Thus apparent changes in a fitted loss-aversion coefficient can reflect changes in noise or gain--loss weight unless their scales are fixed.

At $\eta=0$, $\lambda$ has no effect at all. At $\lambda=1$, the reference enters utility only as the additive constant $-\eta r$ and cancels from within-block comparisons. If all observed $x,y$ lie above $r$, neither the loss slope nor the reference location is identified. Furthermore,
\[
 \frac{\partial\lambda}{\partial m_-}=\frac{s}{s m_+-1}=\frac{s}{\eta}
\]
in (3), displaying weak identification as the gain--loss weight approaches zero. Precise choice probabilities need not produce a precise estimate of $\lambda$ in that region.

Identifying a reference location in each block identifies $\gamma$ only if distinct admissible parameter values give distinct reference maps on the observed support. Randomization of $Z$ supplies variation, but not this injectivity condition. Nor does it establish exclusion: changing a status quo can change information, transaction costs, or perceived quality in addition to the reference.

\section*{A distribution of references is not its mean}
Suppose $\E|R|<\infty$ and the person integrates the random reference $R$ before the choice shock is realized, so the following evaluated utility is finite:
\[
 \overline U(x)=x+\eta\E[v_\lambda(x-R)].                    \tag{4}
\]
This assumption differs from mixing logit probabilities over unobserved realized references. With a continuous reference distribution, bounded differentiation gives
\[
 \overline U'(x)=1+\eta\{1+(\lambda-1)[1-F_R(x)]\}.          \tag{5}
\]
If a density exists, $\overline U''(x)=-\eta(\lambda-1)f_R(x)$. Hence dispersed references can smooth a sharp individual kink into apparent curvature. When $\eta(\lambda-1)>0$ is known, the distribution can formally be recovered from a calibrated derivative curve as
\[
 F_R(x)=1-\frac{\overline U'(x)-1-\eta}{\eta(\lambda-1)}.
\]
An unknown baseline curvature would confound this recovery. At atoms, use one-sided derivatives; the slope jumps encode their masses.

For a direct check, let $R=-1$ or $1$ with equal probability. Its mean is zero, but at $x=0$,
\[
 \E v_\lambda(-R)=\frac{1-\lambda}{2}\ne v_\lambda(0)=0
\]
when $\lambda>1$. This is not merely an irrelevant utility-level shift: between $x=0$ and $x=1/2$, expected gain--loss utility increases by $(1+\lambda)/4$, whereas replacing the reference by its mean gives $1/2$. With $\lambda=2$, these increments are $3/4$ and $1/2$. Thus the catalogue's random-reference variant requires the distributional integral or a justified special case permitting reduction to a point reference.

\section*{Sharp stability bounds from individual parameter intervals}
Suppose there are $D\geq2$ domains, and for each sampled individual there are simultaneous intervals
$\lambda_{i,d}\in[l_{i,d},u_{i,d}]$. They may come from a maintained stochastic-choice model or from partial identification. Define
\[
 L_i=\max\{0,\max_d l_{i,d}-\min_d u_{i,d}\},\qquad
 U_i=\max_{d\ne d'}(u_{i,d}-l_{i,d'})_+,
\]
where $x_+=\max(x,0)$. The maximum pairwise distance equals the range of the domain parameters.
\begin{proposition}
Over the Cartesian product of these intervals, the smallest possible parameter range is $L_i$ and the largest is $U_i$. Consequently the fraction stable within $\varepsilon\geq0$ among $n$ sampled individuals lies in the sharp interval
\[
 \left[\frac1n\sum_i\mathbf1\{U_i\leq\varepsilon\},\quad
       \frac1n\sum_i\mathbf1\{L_i\leq\varepsilon\}\right].     \tag{6}
\]
\end{proposition}
\begin{proof}
Any chosen parameters include one at least $\max_d l_{i,d}$ and one at most $\min_d u_{i,d}$, proving the lower bound on range. If the intervals share a point, choose it in every domain and obtain zero. Otherwise write $a=\min_d u_{i,d}<b=\max_d l_{i,d}$. Every interval intersects $[a,b]$, and selecting in these intersections attains range $b-a$: a domain with upper endpoint $a$ and one with lower endpoint $b$ force the endpoints. These are distinct domains because each interval has ordered endpoints.

Any difference between parameters from distinct domains is bounded above by the corresponding upper-minus-lower endpoints. Selecting endpoints in a maximizing pair attains $U_i$. Therefore every admissible vector is stable exactly when $U_i\leq\varepsilon$, and at least one is stable exactly when $L_i\leq\varepsilon$. When $L_i\leq\varepsilon<U_i$, both stable and unstable vectors exist. Choices can be made independently across individuals, attaining both fractions in (6).
\end{proof}
The distinct-domain restriction in $U_i$ matters: using $\max_d u_{i,d}-\min_d l_{i,d}$ can overstate the largest feasible range when both extremes belong to the same domain. For instance, intervals $[1,10]$ and $[5,6]$ have maximum range five, not nine. Extra joint restrictions across domains or people can narrow (6); the bounds remain valid outer bounds but need not then be sharp.

To translate to a population claim, assume individuals are sampled independently and identically, and the simultaneous interval coverage event has probability at least $1-\delta_m$. For a chosen $\delta_s>0$, set
\[
 h_n=\sqrt{\frac{\log(2/\delta_s)}{2n}}.
\]
Hoeffding's inequality bounds the difference between the empirical fraction of truly stable individuals and its population probability by $h_n$ with probability at least $1-\delta_s$. Combining this event with interval coverage gives the population interval
\[
 \left[\max\{0,\widehat p_L-h_n\},\ \min\{1,\widehat p_U+h_n\}\right] \tag{7}
\]
with coverage at least $1-\delta_m-\delta_s$, where the endpoints in (6) are $\widehat p_L,\widehat p_U$. No independence between coverage and sampling events is needed for the union bound. A lower endpoint at least $1-\alpha$ supports the specified stability criterion under these assumptions. The population exception allowance $\alpha$ is different from the statistical error probabilities $\delta_m,\delta_s$.

\section{Completion Estimate}
50\%

This is a post hoc, heuristic estimate for the full catalogue target.
The attempt derives conditional kink and random-reference identification, exhibits noise-scale and weak-identification failures and gives sharp interval-based bounds for cross-domain stability. Joint individual parameter and reference-process identification from finite longitudinal experiments still requires credible exclusions, comparable scales and empirical coverage procedures for the needed intervals.

\section*{Sources and unresolved identification}
The inspected primary paper Tversky and Kahneman, \emph{Loss Aversion in Riskless Choice: A Reference-Dependent Model} (1991), motivates separating reference dependence from ordinary wealth curvature:
\url{https://www.sscnet.ucla.edu/polisci/faculty/chwe/austen/tversky1991.pdf}.
The author-hosted bibliographic record for K\H{o}szegi and Rabin, \emph{A Model of Reference-Dependent Preferences} (2006), is
\url{https://rabin.scholars.harvard.edu/publications/model-reference-dependent-preferences}.
The calculations above concern the explicitly stated catalogue model and do not claim to establish every feature of either source's framework.

The remaining obstacle is credible joint identification across domains: finitely many binary choices do not by themselves recover an unrestricted distribution of individual parameters and reference processes. Equations (3)--(7) explain which normalizations, experimental variation, and uncertainty bounds would support a stability claim. They provide no empirical evidence that the claim is true, and the full problem remains unresolved.

\end{document}
